\documentclass[11pt]{article}

% Basic packages
\usepackage[margin=1in]{geometry}
\usepackage{amsmath, amssymb, amsfonts}
\usepackage{graphicx}
\usepackage{booktabs}
\usepackage{multirow}
\usepackage{hyperref}
\usepackage{url}
\usepackage[numbers,sort,compress]{natbib}
\usepackage{bm}
\usepackage{mathtools}
\usepackage{xcolor}
\usepackage{algorithm}
\usepackage{algorithmic}
\usepackage{caption}
\usepackage{float}
\graphicspath{{./}{../paperwriting/}{../figures/}}

% --- Hyperlink Settings ---
\hypersetup{
    colorlinks=true,
    linkcolor=blue,
    filecolor=magenta,
    urlcolor=cyan,
    citecolor=blue,
}

\title{\textbf{PocketGNN: A Cross-Modal Framework Unifying Local 3D Pocket Geometry and Global Sequence Semantics for Enzyme Kinetic Prediction}}

\author{
  Zihao Li$^{1}$ \quad Diannan Lu$^{1}$\\[3pt]
  $^{1}$Department of Chemical Engineering, Tsinghua University, Beijing, China\\
  \texttt{l-zh21@mails.tsinghua.edu.cn, ludiannan@tsinghua.edu.cn}
}

\date{}

\begin{document}

\maketitle

\begin{abstract}
The enzyme turnover number ($k_{\text{cat}}$) is a pivotal kinetic parameter for understanding biocatalytic efficiency, yet its accurate prediction remains a grand challenge due to the complex interplay between local physicochemical constraints and global evolutionary context. Existing methods typically bifurcate into sequence-based approaches, which capture evolutionary semantics but miss fine-grained spatial details, or structure-based models, which often suffer from noise in whole-protein representations or lack global context.

To bridge this gap, we propose \textbf{PocketGNN}, a \emph{cross-modal} deep learning framework that synergizes the precision of local 3D geometry with the breadth of global 1D sequence semantics. PocketGNN introduces a high-fidelity graph representation of the active pocket, enriched with a 24-dimensional geometric edge encoding (RBF distances, bond-angle summary statistics, and dihedral-derived summary statistics) to capture the stereochemical determinants of catalysis. Crucially, this local structural view is fused with global evolutionary information extracted from pre-trained protein language models (ESM-2), creating a unified representation that spans spatial scales.

Evaluated on a dataset derived from IntEnzyDB~\citep{yan2022intenzydb}, PocketGNN achieves strong correlation under a standard random split ($r=0.98$, $R^2=0.918$), which we treat as an in-distribution upper bound because homologous enzymes can still overlap across splits. Under a stricter 40\% sequence-identity split designed to test generalization to unseen enzyme families, the model attains $r=0.716$, $R^2=0.387$, and MAE$=0.90$ on an 898-sample test set. To ensure fair baseline comparison, we additionally performed paired evaluations on shared benchmark subsets, where PocketGNN outperforms both CatPred and CataPro across Pearson correlation, $R^2$, MAE, and RMSE (paired permutation test, $p=0.0002$ in both comparisons). Follow-up analyses further show heterogeneous behavior across EC classes and indicate that DiffDock confidence is not a reliable downstream error proxy among successfully constructed pocket graphs. More broadly, PocketGNN reframes enzyme kinetic prediction from family-prior-dominated inference toward scalable, structure-aware representation learning under leakage-aware evaluation, while acknowledging that quantitative mechanistic attribution for mixed categorical and continuous graph features requires additional normalization-aware validation.
\end{abstract}

\paragraph{Keywords:} Enzyme kinetics, Graph neural networks, Molecular docking, Interpretability, Catalytic turnover number, Protein language models

\section{Introduction}
\label{sec:intro}

Enzymes are the orchestrators of life, driving metabolic flux with exquisite specificity. The catalytic turnover number ($k_{\text{cat}}$) serves as the quantitative benchmark of this efficiency~\citep{srinivasan2023guide}. In the era of synthetic biology, the ability to predict $k_{\text{cat}}$ in silico is a holy grail for rational enzyme design and metabolic pathway engineering~\citep{yan2022intenzydb, yu2023unikp}. However, $k_{\text{cat}}$ is an emergent property governed by factors across multiple scales: from the sub-angstrom precision of transition state stabilization in the active site to the global protein dynamics and stability encoded in the entire sequence.


\paragraph{Related Work.}
Predicting enzyme kinetic parameters from sequence and structural information is challenging due to sparse measurements, heterogeneous assay conditions, and substantial experimental noise. Early deep-learning efforts such as \textbf{DLKcat}~\citep{dlkcat} combined protein sequence features with small-molecule representations to predict $k_{\text{cat}}$ at scale, demonstrating that learned representations can capture enzyme–substrate-dependent trends. More recent works increasingly emphasize robust evaluation and broader kinetic coverage: \textbf{UniKP}~\citep{yu2023unikp} proposed a unified framework leveraging pretrained protein language models (pLMs) to predict multiple kinetic parameters. In parallel, \textbf{CatPred}~\citep{catpred} and \textbf{CataPro}~\citep{catapro} introduced comprehensive frameworks incorporating uncertainty estimates and broader parameterization, illustrating the field-wide trend toward multimodal feature integration.

A major driver of progress has been the emergence of large-scale protein language models like \textbf{ProtTrans}~\citep{prottrans} and \textbf{ESM-2/ESMFold}~\citep{esmfold}. These models demonstrate evolutionary-scale modeling and strong representational power, motivating the use of frozen pLM embeddings as global context features. Beyond sequence-only models, structure-aware architectures like \textbf{Uni-Mol}~\citep{unimol} and \textbf{PiFold}~\citep{pifold} have leveraged graph-based encoders for molecular representation and inverse folding. Recent advances such as \textbf{EZSpecificity}~\citep{ezspecificity} employ cross-attention--empowered SE(3)-equivariant GNNs, demonstrating that explicitly modeling interactions between modalities can improve performance.

\paragraph{Our Solution: PocketGNN.}
These developments suggest a converging direction: combining local 3D pocket geometry with global sequence-level evolutionary context. Therefore, we propose \textbf{PocketGNN}, a cross-modal framework designed to unify these complementary views. We posit that accurate kinetic prediction requires a "dual-lens" approach:
\begin{enumerate}
    \item \textbf{Local High-Fidelity Geometry}: We model the active pocket not just as a bag of atoms, but as a geometrically constrained system. By incorporating bond angles and dihedral angles into the graph message passing, we explicitly capture the "short-range" physical constraints essential for catalysis.
    \item \textbf{Global Evolutionary Semantics}: We acknowledge that the pocket does not exist in isolation. By fusing embeddings from ESM-2, we inject "long-range" evolutionary context (e.g., EC family priors, conservation patterns) into the prediction.
\end{enumerate}

This fusion of \emph{local 3D physics} and \emph{global 1D evolution} allows PocketGNN to outperform existing baselines, providing a robust tool for functional enzyme mining.

Beyond predictive accuracy, we view this work as addressing a central issue in enzyme machine learning: many reported gains can still be dominated by enzyme-family priors or optimistic data splits, whereas catalytic turnover is fundamentally constrained by the local stereochemical organization of the active site. Our goal is therefore not only to build a stronger predictor, but also to test whether explicitly modeling catalytic pocket geometry remains useful under leakage-aware homology splits that better reflect realistic deployment on unseen enzyme families.

\paragraph{Main Contributions.}
The main contributions of this work are fourfold:
\begin{enumerate}
    \item We introduce \textbf{PocketGNN}, a cross-modal framework that combines local 3D pocket geometry with global sequence semantics for enzyme $k_{\text{cat}}$ prediction.
    \item We evaluate the model under a strict \textbf{40\% homology-aware split} and \textbf{paired shared-sample baseline comparisons}, providing a more leakage-aware assessment than direct comparison of metrics computed on mismatched subsets.
    \item We analyze component sensitivity, feature-scale effects, and failure modes across EC classes and docking-confidence strata, clarifying both the strengths and current interpretability limits of catalytic pocket representations.
    \item We demonstrate that an \textbf{automated docking-to-pocket pipeline} can serve as a scalable preprocessing strategy for structure-aware enzyme kinetic modeling without manual binding-site annotation.
\end{enumerate}

\section{Methods}
\label{sec:methods}

\subsection{Data Curation and Pipeline}
\label{sec:data_pipeline}

We curate our dataset from IntEnzyDB~\citep{yan2022intenzydb}, filtering for entries with valid UniProt IDs, substrate SMILES, and experimental $k_{\text{cat}}$ values. To ensure data quality, we exclude extreme outliers and select entries where $k_{\text{cat}} \in [10^{-2}, 10^{6}] \text{ s}^{-1}$. Our processed dataset contains 9,059 enzyme--substrate pairs with experimentally measured $k_{\text{cat}}$ values, covering 6 major EC classes (EC 1--6).

The data processing pipeline (Figure \ref{fig:framework}) proceeds as follows:
\begin{enumerate}
    \item \textbf{Structure Acquisition}: Enzyme structures are sourced from the PDB or predicted using ESMFold~\citep{esmfold}. Substrate 3D conformers are generated from SMILES using RDKit with UFF energy minimization.
    \item \textbf{Pocket Extraction via DiffDock}: We perform molecular docking using DiffDock~\citep{diffdock2023}, a score-based diffusion model that predicts binding poses from protein structure and ligand SMILES alone, \emph{without requiring manual specification of a binding site}. The active pocket is then defined as all protein atoms within 5.0~\AA{} of the top-ranked docked pose.
    \item \textbf{Graph Construction}: The pocket is converted into a radius graph ($r_{\text{cut}} = 4.0 \text{ \AA}$) for input into the GNN.
\end{enumerate}
For homology-aware evaluation, we cluster enzyme sequences with MMseqs2 at 40\% identity, obtaining 1,530 clusters in total, and split clusters rather than individual samples to avoid family-level leakage.

\paragraph{Why DiffDock over AutoDock Vina.}
Traditional physics-based docking methods such as AutoDock Vina~\citep{trott2010autodock} require explicit specification of the binding-site search box (center coordinates and dimensions), which is labor-intensive and impractical for large-scale datasets with no prior structural knowledge. Critically, when protein structures are predicted from sequences (e.g., via ESMFold) rather than derived from crystal structures, no reference ligand coordinates exist to guide search box placement. DiffDock eliminates this requirement through an SE(3)-invariant diffusion process that automatically normalizes protein coordinates to their centroid and searches the entire protein surface, making it robust to arbitrary input orientations.

In the original DiffDock benchmark on the full Astex Diverse Set ($n$=85), DiffDock achieves a top-1 success rate of 38.2\% (RMSD $<$ 2~\AA) in a \emph{blind docking} setting where no binding site is provided, compared to 20.9\% for AutoDock Vina with an explicit search box~\citep{diffdock2023}. To validate pose accuracy on our pipeline, we performed a small-scale re-docking experiment on 5 representative complexes (Table~\ref{tab:docking_rmsd}). In this \emph{re-docking} setting (where native ligand coordinates guide Vina's search box), both methods achieve reasonable accuracy: Vina attains 100\% success with mean RMSD 1.31$\pm$0.61~\AA, while DiffDock achieves 80\% success with mean RMSD 1.69$\pm$1.19~\AA. However, this comparison is inherently biased toward Vina, as it assumes \emph{a priori} knowledge of the binding site---precisely the information unavailable in our sequence-to-structure prediction pipeline. For processing thousands of enzyme--substrate pairs from predicted structures without manual annotation, DiffDock's \textbf{fully automated blind docking} capability is essential.

\begin{table}[h]
\centering
\caption{Re-docking validation on 5 Astex Diverse Set complexes. Both methods achieve reasonable accuracy in this controlled setting; DiffDock's advantage is automation (no binding-site specification required).}
\label{tab:docking_rmsd}
\begin{tabular}{lcccc}
\toprule
Method & Binding Site & Success Rate & Mean RMSD & Best RMSD \\
\midrule
AutoDock Vina~\citep{trott2010autodock} & Required & 100\% (5/5) & 1.31$\pm$0.61\,\AA & 0.30\,\AA \\
DiffDock~\citep{diffdock2023} & \textbf{Not needed} & 80\% (4/5) & 1.69$\pm$1.19\,\AA & \textbf{0.38\,\AA} \\
\bottomrule
\end{tabular}
\end{table}

\begin{figure}[htbp]
    \centering
    \includegraphics[width=\textwidth]{Architecture_EN.jpg}
    \caption{The cross-modal framework of PocketGNN. The model integrates a high-fidelity local graph encoder (capturing distance, angle-derived, and torsion-derived edge statistics) with a global sequence encoder (ESM-2), bridging the gap between fine-grained geometry and evolutionary semantics.}
    \label{fig:framework}
\end{figure}

\subsection{PocketGNN Architecture}
\label{sec:model}

PocketGNN employs a multi-branch architecture to fuse information across modalities.

\subsubsection{Local Branch: Geometry-Enhanced Graph Encoder}
The active site is modeled as a geometric graph.
\begin{itemize}
    \item \textbf{Node Features (52-dim)}: Atom-level descriptors including element type, residue identity, ligand/protein indicator, nearest-neighbor distance, electronic features, basic atomic properties, secondary-structure one-hot features, and scaled solvent accessibility.
    \item \textbf{Geometric Edge Encoding (24-dim)}: To capture the stereochemistry, we employ an edge featurization $\mathbf{e}_{ij}$ comprising:
    \begin{itemize}
        \item \textbf{Radial}: 16-dim Gaussian RBF encoding of distances.
        \item \textbf{Angular}: 4-dim summary statistics derived from bond-angle neighborhoods around each edge.
        \item \textbf{Torsional}: 4-dim summary statistics derived from local dihedral configurations, capturing conformational flexibility.
    \end{itemize}
\end{itemize}
A Graph Attention Network (GAT)~\citep{velickovic2018graph} processes this graph, where attention weights are modulated by these geometric edge features, ensuring that the message passing respects the spatial configuration of the pocket.

\subsubsection{Global Branch: Sequence Semantics Fusion}
To incorporate long-range dependencies, we utilize ESM-2 embeddings. For each enzyme, a global sequence vector $\mathbf{s}_{\text{seq}}$ is extracted. This vector encapsulates the evolutionary history and global structural propensities of the protein.

\subsubsection{Cross-Modal Late Fusion}
The local structural representation $\mathbf{h}_{\text{pocket}}$ (from graph pooling) and the global sequence representation $\mathbf{s}_{\text{proj}}$ (from ESM-2 projection) are concatenated:
\begin{equation}
    \mathbf{z}_{\text{final}} = \text{MLP}([\mathbf{h}_{\text{pocket}} \parallel \mathbf{s}_{\text{proj}}])
\end{equation}
To avoid dimensional imbalance between the high-dimensional sequence embedding (640-1280 dim) and the graph representation (128-256 dim), we employ a projection layer that maps the ESM-2 embedding to the same dimension as the graph features before concatenation. This late fusion allows the model to dynamically weigh the immediate physical constraints of the pocket against the global stability implied by the sequence.

\section{Experiments}
\label{sec:experiments}

\subsection{Experimental Setup}
We report two evaluation settings. (i) \textbf{Random split}: 80\% training and 20\% validation for an in-distribution upper bound. (ii) \textbf{Homology-aware split}: we cluster sequences with MMseqs2 at 40\% identity and construct train/val/test splits at the cluster level, ensuring no test sequence exceeds 40\% identity to any training sequence. This split serves as the primary indicator of out-of-distribution generalization. Throughout the paper we emphasize Pearson correlation and $R^2$, and additionally report MAE and RMSE for the homology-split and shared-sample comparisons.

\subsection{Data Quality and Feature Analysis}
\label{sec:data_analysis}

To ensure that our graph representations capture meaningful physicochemical signals rather than spurious correlations, we performed a comprehensive data-level diagnosis. Pearson correlation analysis revealed that several geometric and electronic features exhibit significant correlation ($|r| > 0.3$) with catalytic rates. Furthermore, a Random Forest regressor trained solely on aggregated statistical descriptors achieved a test Pearson correlation of $\sim0.55$. While this suggests that simple descriptors alone cannot fully solve the problem, it confirms the existence of \emph{nonlinear but extractable} signal within the pocket geometry, validating our focus on the active site.

\subsection{Performance Comparison}

PocketGNN demonstrates superior performance, effectively leveraging its cross-modal design.

\paragraph{In-Distribution Performance (Random Split).}
As shown in Table \ref{tab:results}, the full model achieves a Pearson $r$ of 0.98 and $R^2$ of 0.918 for $\log_{10}(k_{cat})$. The improvement over the structure-only and sequence-only baselines highlights the synergy of combining local geometry with global semantics. We note that this setting can overestimate performance due to family-level overlap, so we treat the homology-aware split as the more faithful measure of generalization.

\begin{table}[h]
\centering
\caption{Quantitative prediction results for $\log_{10}(k_{\text{cat}})$ on the random-split validation set.}
\label{tab:results}
\begin{tabular}{lccc}
\toprule
Model & MSE & $R^2$ & Pearson $r$ \\
\midrule
Sequence Baseline & 2.227 & -30.15 & 0.18 \\
Sequence Enhanced & 0.28 & -0.32 & 0.82 \\
PGNN-Basic (Structure Only) & 0.79 & 0.665 & 0.95 \\
\textbf{PocketGNN (Cross-Modal)} & \textbf{0.19} & \textbf{0.918} & \textbf{0.98} \\
\bottomrule
\end{tabular}
\end{table}

\paragraph{Out-of-Distribution Generalization (Homology Split).}
To assess robustness against data leakage, we evaluated the model on the strict 40\% sequence identity split (898 test samples). In this challenging setting, where no test enzyme shares $>$40\% sequence identity with any training enzyme, PocketGNN achieves a Pearson correlation of \textbf{0.716} (with $R^2=0.387$, MAE$=0.90$). The moderate $R^2$ reflects the inherent difficulty of predicting kinetic rates for unseen enzyme families, while the relatively strong linear correlation indicates that the model still captures transferable trends under family-level distribution shift. In this setting, the automated DiffDock-based pocket extraction pipeline serves as a scalable structure preprocessing step, making it feasible to build pocket graphs without manually curated binding sites.

\paragraph{Comparison with State-of-the-Art Methods.}
Direct comparison between public baselines is complicated by non-identical sample coverage. We therefore evaluated all methods on \emph{shared sample subsets} and performed paired statistical tests (Table \ref{tab:sota_comparison}). On the 3,439 samples shared with CatPred, PocketGNN improves Pearson correlation from 0.608 to 0.864 and reduces MAE from 0.942 to 0.479. On the 1,455 samples shared with CataPro, PocketGNN improves Pearson correlation from 0.683 to 0.868 and reduces MAE from 0.892 to 0.487. Paired permutation tests confirm that these gains are significant for Pearson correlation, $R^2$, MAE, and RMSE in both comparisons ($p=0.0002$ in each comparison). This shared-sample analysis is more rigorous than comparing metrics computed on different subsets and shows that PocketGNN's cross-modal representation provides a consistent advantage over both sequence-dominant and alternative multimodal baselines.

\begin{table}[h]
\centering
\caption{Paired comparison with public baselines on shared benchmark subsets. Each two-row block uses the exact same samples for both methods; paired permutation tests indicate significant improvements for PocketGNN across Pearson correlation, $R^2$, MAE, and RMSE in both blocks ($p=0.0002$ for PocketGNN vs.\ CatPred and PocketGNN vs.\ CataPro).}
\label{tab:sota_comparison}
\begin{tabular}{llccccc}
\toprule
Shared subset & Method & N & Pearson $r$ & $R^2$ & MAE & RMSE \\
\midrule
CatPred overlap & CatPred (2025)~\citep{catpred} & 3439 & 0.608 & 0.165 & 0.942 & 1.405 \\
CatPred overlap & \textbf{PocketGNN (ours)} & 3439 & \textbf{0.864} & \textbf{0.745} & \textbf{0.479} & \textbf{0.775} \\
\midrule
CataPro overlap & CataPro (2025)~\citep{catapro} & 1455 & 0.683 & 0.444 & 0.892 & 1.178 \\
CataPro overlap & \textbf{PocketGNN (ours)} & 1455 & \textbf{0.868} & \textbf{0.752} & \textbf{0.487} & \textbf{0.787} \\
\bottomrule
\end{tabular}
\end{table}

Figure~\ref{fig:scatter} visualizes the correlation between predicted and ground-truth values on the homology-split test set ($r=0.716$), showing coherent clustering along the diagonal despite the challenging OOD setting.

As a calibration diagnostic, we also binned the 898 homology-split test predictions by true $\log_{10}(k_{\text{cat}})$ values. The model shows a regression-to-the-mean pattern: very low-rate enzymes tend to be over-predicted, whereas very high-rate enzymes tend to be under-predicted. This explains why the model can preserve a useful rank/trend signal while still showing moderate $R^2$ and large errors in the extreme kinetic ranges.

To test whether these gains hold across sample strata, we further partitioned the shared-sample comparisons by enzyme sequence length and ligand heavy-atom count. PocketGNN remained significantly better than CatPred in every tertile (Pearson advantage 0.207--0.286; MAE advantage 0.391--0.547) and also consistently outperformed CataPro in every tertile (Pearson advantage 0.158--0.227; MAE advantage 0.355--0.433), with all within-bin permutation tests yielding $p=0.0005$. These stratified results indicate that the improvement is not confined to one narrow region of sequence or substrate complexity.

\begin{figure}[htbp]
    \centering
    \includegraphics[width=0.65\textwidth]{test_kcat_prediction_scatter.png}
    \caption{True vs.\ predicted $\log_{10}(k_{\text{cat}})$ on the 40\% homology-split test set (898 samples). Pearson $r=0.716$, $R^2=0.387$, MAE$=$0.90. Each point is one enzyme--substrate pair; color indicates local density.}
    \label{fig:scatter}
\end{figure}

\begin{figure}[htbp]
    \centering
    \includegraphics[width=0.85\textwidth]{diffdock_vs_vina_comparison.png}
    \caption{Re-docking validation on 5 Astex Diverse Set complexes. \textbf{Left}: Per-complex RMSD comparison. \textbf{Right}: Summary statistics. Both methods achieve reasonable accuracy (all RMSD $<$4\,\AA); Vina shows higher consistency while DiffDock achieves the lowest individual RMSD. The key advantage of DiffDock is automation---no binding-site specification required.}
    \label{fig:docking_comparison}
\end{figure}

Because PocketGNN depends on docking-derived pockets, we also assessed whether DiffDock confidence can serve as a proxy for downstream prediction error. Across 940 successfully docked samples, DiffDock confidence shows no significant correlation with absolute prediction error (Pearson $r=0.027$, $p=0.41$; Spearman $\rho=0.036$, $p=0.27$). Performance also remains broadly similar across confidence tertiles (Pearson $r=0.834$--0.886; MAE$=0.449$--0.491). We therefore interpret DiffDock confidence as an unreliable downstream error proxy in this setting, not as a substitute for direct pose-quality validation; dense regions of the confidence scatter naturally contain more visible large-error points simply because they contain more samples.

\subsection{Component and Feature-Scale Diagnostics}
\label{sec:ablation}

The graph representation combines binary categorical indicators, count-like electronic descriptors, elemental physicochemical properties, radial basis encodings, and bounded angular summaries. Because these feature groups live on different numerical scales, we do not use raw gradient magnitude as a definitive quantitative measure of mechanistic importance. Instead, we treat feature-level attribution as exploratory and focus the main manuscript on performance, shared-sample benchmarking, EC-wise heterogeneity, and docking-robustness diagnostics. Future component studies should use train-set-fitted standardization and group perturbation experiments so that distance, angular, categorical, and continuous channels can be compared under a common attribution protocol.

\subsection{EC-wise Performance Analysis}
\label{sec:ec_analysis}

To further understand model behavior across enzyme classes, we analyzed performance metrics by EC number on the 40\% homology test set using the same prediction file as the main evaluation. This analysis covers all 898 test samples with EC labels in the processed dataset. Table~\ref{tab:ec_wise} shows that behavior is heterogeneous across EC classes: EC 2 and EC 6 show relatively favorable performance, whereas EC 5 maintains high correlation but has weak $R^2$ and high MAE. We interpret this as descriptive evidence of EC-dependent calibration and sample-size effects, not as proof that all enzyme classes generalize equally well.

\begin{table}[h]
\centering
\caption{Performance by EC class on the 40\% homology test set, aligned to the main homology-split prediction file. EC 7 contains only one sample and is omitted from the table.}
\label{tab:ec_wise}
\begin{tabular}{lccccc}
\toprule
EC Class & Name & N & Pearson $r$ & $R^2$ & MAE \\
\midrule
EC 1 & Oxidoreductases & 703 & 0.667 & 0.350 & 0.900 \\
EC 2 & Transferases & 67 & 0.773 & 0.503 & 0.559 \\
EC 3 & Hydrolases & 36 & 0.885 & 0.376 & 0.932 \\
EC 4 & Lyases & 42 & 0.813 & 0.422 & 1.047 \\
EC 5 & Isomerases & 31 & 0.843 & 0.059 & 1.383 \\
EC 6 & Ligases & 18 & 0.973 & 0.615 & 0.932 \\
\midrule
\textbf{Overall} & & \textbf{898} & \textbf{0.716} & \textbf{0.387} & \textbf{0.901} \\
\bottomrule
\end{tabular}
\end{table}

\begin{figure}[htbp]
    \centering
    \includegraphics[width=0.85\textwidth]{../results/ec_analysis_corrected/fig_ec_wise_performance_corrected.png}
    \caption{EC-wise performance on the 40\% homology-split test set using predictions aligned to the main evaluation. The figure highlights EC-dependent heterogeneity rather than uniformly strong performance across all classes.}
    \label{fig:ec_wise}
\end{figure}

\subsection{Interpretability and Mechanistic Insights}
\label{sec:interpretability}

To understand \emph{why} PocketGNN achieves strong predictive performance, we conducted qualitative saliency analysis on representative cases. We intentionally treat feature-level gradient attributions as exploratory because the graph inputs mix one-hot categorical variables, count-like electronic descriptors, and continuous physicochemical quantities on different numerical scales. Raw zero-baseline Integrated Gradients therefore does not provide a physically meaningful interpolation path for every feature type.

\subsubsection{Exploratory Feature Attribution}

Preliminary attribution experiments suggested that predictions can be sensitive to geometric edge channels, but we do not use the resulting percentages as quantitative mechanistic evidence in this manuscript. A rigorous comparison between distance, angular, categorical, and continuous channels should use train-set-fitted feature standardization together with group masking or permutation tests. We therefore present the current attribution results only as motivation for future component-level validation, while relying on prediction performance and failure-mode analyses for the main claims.

\subsubsection{Case Study: Saliency on Catalytic Residues}

As an illustrative case study, we inspected Input $\times$ Gradient saliency on Human Aldehyde Dehydrogenase 2 (ALDH2, P30838). The visualization (Figure \ref{fig:case_study}) shows that high-saliency regions overlap with the catalytic nucleophile (Cys302), the proton abstractor (Glu268), and the substrate binding cleft. This case supports the chemical plausibility of the learned pocket representation, but it should be interpreted as qualitative evidence rather than a standalone quantitative benchmark of mechanism recovery.

\begin{figure}[htbp]
    \centering
    \includegraphics[width=0.7\textwidth]{visualize.png}
    \caption{Visualization of atomic contributions for ALDH2. High-saliency atoms (red) coincide with known catalytic residues (Cys302, Glu268), providing a qualitative example that the model attends to chemically relevant pocket regions.}
    \label{fig:case_study}
\end{figure}

\subsection{Code and Data Availability}
All code, trained models, and processed datasets are available at \url{https://github.com/RektonLee/PGNN_final}. The raw data is derived from IntEnzyDB~\citep{yan2022intenzydb}, a publicly available database. Pre-trained models and processed datasets are available upon reasonable request.

\subsection{Limitations and Future Work}
\label{sec:limitations}
Our study has several limitations that warrant future investigation:

\begin{itemize}
    \item \textbf{Experimental Conditions}: Unlike MPEK~\citep{mpek2024}, we do not explicitly model temperature, pH, or organism-specific effects. This may limit accuracy when predicting kinetics under conditions different from training data.

    \item \textbf{Uncertainty Quantification}: Unlike CatPred~\citep{catpred}, we do not provide prediction confidence intervals. Future work will incorporate quantile regression or ensemble-based uncertainty estimation.

    \item \textbf{Multi-substrate Reactions}: Our current framework handles single-substrate reactions. Extending to multi-substrate systems (e.g., reactant concatenation as in CatPred) would broaden applicability.

    \item \textbf{Run-to-run Variability}: The main tables report selected checkpoints for each split rather than mean$\pm$std across multiple random seeds. Future work should add multi-seed evaluation to quantify training variance more rigorously.

    \item \textbf{Docking Dependency}: The quality of pocket extraction depends on DiffDock's accuracy. Although DiffDock confidence is not a reliable proxy for downstream prediction error on 940 successfully constructed graphs, confidence scores are not equivalent to experimental pose accuracy. Future work should expand validation on holo structures with known ligand poses.

    \item \textbf{Feature Scaling and Attribution}: The current graph representation mixes one-hot categorical variables, count-like electronic descriptors, and continuous physicochemical features. Training and inference use the same feature scales, so the predictive results remain internally consistent, but raw gradient attribution can be scale-dependent. Future work should use train-set-fitted feature standardization and group perturbation to re-evaluate attribution stability.

    \item \textbf{Data Quality}: Our study focuses on a curated high-quality dataset, and performance on low-quality or incomplete data remains to be tested.
\end{itemize}

While the homology split demonstrates generalization capability, the drop in $R^2$ indicates room for improvement in transfer learning across disparate enzyme families. Future work could explore pre-training strategies specifically for pocket geometries to further bridge this gap.

\section{Conclusion}
\label{sec:conclusion}

PocketGNN represents a practical step forward in \emph{cross-modal} enzyme modeling. By fusing the "short-range" precision of geometry-enhanced pocket graphs with the "long-range" breadth of protein language models, we achieve strong kinetic prediction performance ($r=0.98$ under random splitting and $r=0.716$ under strict homology splitting). Shared-sample paired comparisons further show that PocketGNN outperforms both CatPred and CataPro on benchmark subsets, while docking robustness analysis indicates that DiffDock confidence is not a reliable downstream error proxy once valid pocket graphs have been constructed. EC-wise analysis further reveals heterogeneous behavior across enzyme classes, emphasizing both the promise and the remaining calibration challenges of structure-aware $k_{\text{cat}}$ prediction.

Taken together, these results suggest a broader message for enzyme informatics: accurate kinetic prediction should move beyond family-level memorization and toward \emph{catalysis-aware} representations that explicitly encode the geometric constraints of the active pocket, and such structure-aware modeling can be scaled through automated docking-based preprocessing. We therefore view PocketGNN not only as a predictive model, but also as evidence that leakage-aware evaluation plus local stereochemical modeling is a productive direction for future enzyme foundation models and functional screening pipelines.

% =======================================================
% ACKNOWLEDGMENTS
% =======================================================
\section*{Acknowledgments}
This work was supported by the National Key Research and Development Program of China and the National Natural Science Foundation of China. We thank the Tsinghua University High-Performance Computing Center for providing computational resources. We acknowledge IntEnzyDB for providing the curated enzyme kinetics database and the developers of DiffDock, ESM-2, and PyTorch Geometric for making their tools publicly available. Z.L. acknowledges helpful discussions with colleagues in the Chemical Engineering Department.

\section*{Author Contributions}
Z.L. conceived the project, developed the methodology, performed all experiments, and wrote the manuscript. D.L. supervised the research and provided critical feedback on the manuscript. All authors reviewed and approved the final version.

\section*{Competing Interests}
The authors declare no competing financial interests or personal relationships that could have appeared to influence the work reported in this paper.

% =======================================================
% REFERENCES
% =======================================================
\begin{thebibliography}{16}

\bibitem{srinivasan2023guide}
Srinivasan B. A guide to enzyme kinetics in early drug discovery. \textit{FEBS Journal}, 2023, 290(9): 2292-2305.

\bibitem{dlkcat}
Li F, Li C, Wang M, et al. DLKcat: deep learning for kcat prediction using protein sequence and substrate structure. \textit{Bioinformatics}, 2022, 38(16): 3904-3910.

\bibitem{yu2023unikp}
Yu H, Deng H, He J, et al. UniKP: a unified framework for the prediction of enzyme kinetic parameters. \textit{Nature Communications}, 2023, 14(1): 8211.

\bibitem{catpred}
Boorla VS, Maranas CD. CatPred: a comprehensive framework for deep learning in vitro enzyme kinetic parameters. \textit{Nature Communications}, 2025, 16: 2072.

\bibitem{catapro}
Yang K, et al. Robust enzyme discovery and engineering with deep learning using CataPro. \textit{Nature Communications}, 2025, 16: 2736.

\bibitem{prottrans}
Elnaggar A, et al. ProtTrans: Toward Understanding the Language of Life Through Self-Supervised Learning. \textit{IEEE TPAMI}, 2022.

\bibitem{esmfold}
Lin Z, Akin H, Rao R, et al. Evolutionary-scale prediction of atomic-level protein structure with a language model. \textit{Science}, 2023, 379(6637): 1123-1130.

\bibitem{unimol}
Zhou G, et al. Uni-Mol: A Universal 3D Molecular Representation Learning Framework. \textit{ICLR}, 2023.

\bibitem{pifold}
Gao Z, et al. PiFold: Toward effective and efficient protein inverse folding. \textit{ICLR}, 2023.

\bibitem{ezspecificity}
Li X, et al. Enzyme specificity prediction using cross attention graph neural networks. \textit{Nature Communications}, 2025.

\bibitem{yan2022intenzydb}
Yan B, Ran X, Gollu A, et al. IntEnzyDB: an Integrated Structure-Kinetics Enzymology Database. \textit{Journal of Chemical Information and Modeling}, 2022, 62(22): 5841-5848.

\bibitem{kroll2021deep}
Kroll A, Engqvist MKM, Heckmann D, et al. Deep learning allows genome-scale prediction of Michaelis constants from structural features. \textit{PLoS Biology}, 2021, 19(10): e3001402.

\bibitem{trott2010autodock}
Trott O, Olson AJ. AutoDock Vina: improving the speed and accuracy of docking with a new scoring function. \textit{Journal of Computational Chemistry}, 2010, 31(2): 455-461.

\bibitem{velickovic2018graph}
Velickovic P, Cucurull G, Casanova A, et al. Graph Attention Networks. \textit{ICLR}, 2018.

\bibitem{diffdock2023}
Corso G, Stärk H, Jing B, et al. DiffDock: Diffusion Steps, Twists, and Turns for Molecular Docking. \textit{ICLR}, 2023.

\bibitem{mpek2024}
Wang Y, et al. MPEK: a multitask deep learning framework based on pretrained language models for enzymatic reaction kinetic parameters prediction. \textit{Briefings in Bioinformatics}, 2024, 25(5): bbae387.

\bibitem{sundararajan2017axiomatic}
Sundararajan M, Taly A, Yan Q. Axiomatic Attribution for Deep Networks. \textit{Proceedings of the 34th International Conference on Machine Learning (ICML)}, 2017, 70: 3319-3328.

\end{thebibliography}

\end{document}
